\documentclass[11pt,a4paper]{article}
\usepackage[margin=1in]{geometry}
\usepackage{amsmath,amssymb,amsthm}
\usepackage{booktabs}
\usepackage{array}
\usepackage{fontspec}
\setmainfont{lmroman10-regular.otf}[
  BoldFont=lmroman10-bold.otf,
  ItalicFont=lmroman10-italic.otf,
  BoldItalicFont=lmroman10-bolditalic.otf,
]
\setmonofont{Menlo}[Scale=0.86, Ligatures=NoCommon]
\usepackage{microtype}
\usepackage{xcolor}
\definecolor{seclink}{HTML}{1F4E8C}
\usepackage[colorlinks=true,linkcolor=seclink,citecolor=seclink,urlcolor=blue]{hyperref}
\usepackage{parskip}
\usepackage{fancyvrb}
\usepackage[most]{tcolorbox}

\definecolor{yellowbg}{HTML}{FBF3D5}
\definecolor{pinkbg}{HTML}{FADDE1}
\definecolor{bluebg}{HTML}{E8EEF7}

\tcbset{plainbg/.style={boxrule=0pt, sharp corners, colframe=white,
  left=4mm, right=4mm, top=3mm, bottom=3mm, before skip=3mm, after skip=3mm,
  before upper={\setlength{\parskip}{0.6\baselineskip}}}}
\newtcolorbox{sectionbox}{plainbg, colback=yellowbg}
\newtcolorbox{alarmbox}{plainbg, colback=pinkbg}
\newtcolorbox{notebox}{plainbg, colback=bluebg}

\DefineVerbatimEnvironment{term}{Verbatim}{fontsize=\footnotesize, xleftmargin=0mm}
\DefineVerbatimEnvironment{code}{Verbatim}{fontsize=\small, xleftmargin=0mm}

\newtheorem{theorem}{Theorem}
\newtheorem*{claim}{Claim}
\newtheorem*{obs}{Observation}

\title{\vspace{-2em}Almost Computable}
\author{\normalsize The Myhill--Nerode quotient of a directed container\\[0.2em]
\normalsize is the cover its cohomology class demands}
\date{\normalsize Claude, 7 September 2026}

\begin{document}
\maketitle
\vspace{-1em}

\begin{sectionbox}
Two pieces of work in this collection were written two months apart, with no
shared vocabulary and no cross-reference. \emph{Coloured Graphs as Directed
Containers} is from July. It builds the path directed container of a graph, asks
when a labelling of its positions is computable, and grades the failure by how
fast a partition refinement splits. The federation programs are from September.
They put three servers in a ring, read off a class in $H^1(S^1;\mathbb{Z}/2)$,
and pay for a centre by passing to a double cover.

The July document does not contain the words cohomology, sheaf, monodromy,
cocycle or cover. Not once. The September documents do not contain the words
Myhill--Nerode, refinement or computable labelling.

They are the same construction. Running the July refinement on the September
ring produces the September double cover, node for node, with the section as its
state set. This note demonstrates that, states what each side gains from the
other, and corrects one claim in \texttt{sheaf-federation.pdf} that the
identification shows to be wrong.
\end{sectionbox}

\section{Where the files are}

Full paths, verified at the time of writing.

\medskip
\textbf{The July work.} A repository at \texttt{git@github.com:RaggedR/category.git}.

\begin{term}
/Users/robin/git/category/category-theory/knowledge-graph/
    coloured-graphs-as-directed-containers.tex     2072 lines,  78514 bytes
    coloured-graphs-as-directed-containers.pdf                 188224 bytes
    refinement_experiment.py                         54 lines,   1569 bytes
    series_mod2.py                                   62 lines,   1707 bytes
    split_locations.py                               57 lines,   1984 bytes
\end{term}

The relevant sections of that document are \S\,``Computations on directed
graphs'', \S\,``Directed Containers'' with Theorem~\ref{thm:july} below,
\S\,``The computable case: coKleisli collapses to pure functions'' with its two
closing remarks, and \S\,``Almost computable: Sturmian sequences'' with its
experiment.

\medskip
\textbf{A shorter document of the same name}, which is a different text and not
a copy. No git remote.

\begin{term}
/Users/robin/git/kodamai-knowledge-graph/
    coloured-graphs-as-directed-containers.tex      409 lines,  17649 bytes
    coloured-graphs-as-directed-containers.pdf
\end{term}

\medskip
\textbf{Other directed-container material.} Of the two directories below only
\texttt{macbeth-seed} has a git remote, and it is
\texttt{git@github.com:RaggedR/macbeth-seed.git}\@.

\begin{term}
/Users/robin/git/typed-agents/hris-directed-containers.tex      643 lines
/Users/robin/git/typed-agents/directed-graphs.tex              1004 lines
/Users/robin/git/macbeth-seed/learning/directed-containers.tex  964 lines
/Users/robin/git/macbeth-seed/learning/directed-categories.tex
/Users/robin/git/macbeth-seed/learning/directed-containers-tax.tex
/Users/robin/git/macbeth-seed/pdf/directed-containers/
/Users/robin/git/macbeth-seed/notes/DIRECTED_CONTAINER.md
/Users/robin/git/macbeth-seed/notes/LYRA_DIRECTED_CONTAINER.md
/Users/robin/git/macbeth-seed/notes/UMBRAL_DIRECTED_CONTAINER.md
/Users/robin/git/macbeth-seed/notes/GA_DIRECTED_CONTAINER.md
/Users/robin/git/macbeth-seed/notes/ETHEREUM_COMPOSE_DIRECTED_CONTAINER.md
\end{term}

Also relevant, and both with remotes:
\texttt{git@github.com:RaggedR/social-network.git}, which treats social networks
as directed container comonads, and the private \texttt{RaggedR/marie-claire},
which is a Haskell system built on directed containers with a Lean twin.

\medskip
\textbf{The September work.} A repository at
\texttt{git@github.com:RaggedR/federation.git}, public.

\begin{term}
/Users/robin/git/gold/federation/
    federation.pdf            The Missing Trace
    sheaf-federation.pdf      A Sheaf Theory of Federation
    seven-federations.pdf     The Root Outside the Set
    mobius.pdf                the program, in container language
    cocycle.ts                the mathematics, no IO           404 lines
    mobius.server.ts          one arc of the cover
    mobius.ts                 the CLI
    mobius-type-tests.ts      the cohomology, asserted to the compiler
    lib/algebra.ts            the four combinators, first class
    lib/wire.ts               the boundary
\end{term}

\section{The two constructions}

\subsection{July: the path container of a graph, and when a labelling is computable}

Take a directed graph $G$. The container has one shape per vertex, and the
positions over a vertex $v$ are the finite paths out of $v$.

\begin{theorem}[July, \S\,Directed Containers]
\label{thm:july}
Every computable directed graph $G$ determines a computable directed container
$(V_G, P_G)$, with root the empty path, sub-shape $v \downarrow p =
\mathrm{end}(p)$, and shift given by concatenation.
\end{theorem}

Positions are paths rather than edges for a stated reason. The shift has to
compose two positions, and the concatenation of two edges is not an edge, while
the concatenation of two paths is a path.

A directed container is a comonad, so there is a coKleisli category. Label each
vertex with a value in $X$ and each edge with a function of the right type. A
labelling $\ell$ of paths is called \emph{computable} when it is a fold of the
edge functions,
\[
  \ell(p) = (f_{e_n} \circ \cdots \circ f_{e_1})(x_0),
  \qquad
  F(p \oplus q) = F(q) \circ F(p),
  \qquad
  F(\mathrm{root}) = \mathrm{id}.
\]
Those two equations are exactly the directed container laws. The counit reads
the present value. The comultiplication decorates every path with its whole
future. The collapse theorem says that on computable labellings the future is
determined by the present, so every coKleisli map factors through the counit and
coKleisli composition becomes ordinary composition.

\begin{notebox}
\textbf{The sentence this note is built on}, quoted from
\S\,``The computable case''.

\emph{Hidden state is precisely the failure of a map to factor through
$\varepsilon$.}
\end{notebox}

Then an algorithm. Group paths by their end vertex. Refine by comparing
continuations one level deeper. Repeat until a round splits nothing. The
resulting classes are the Myhill--Nerode classes, and their number at a vertex
is the memory an implementation needs there. The document grades the outcome. A
labelling is $N$-not-computable if a split occurs at step $N$, and computable
exactly when that happens for only finitely many $N$. Sturmian labellings split
once per level forever, as slowly as an infinite refinement can, and those are
the almost computable ones.

\subsection{September: three arcs, and the class the ring carries}

Cover the circle with three arcs, each slightly over a third. Coefficients in
$\mathbb{Z}/2$, one sign per overlap. The first cohomology is $\mathbb{Z}/2$,
so there are two classes. Product $+1$ is the cylinder and product $-1$ is the
M\"obius band. Three servers form the ring, each a container and each a morphism,
and the sign is applied in the backward map and nowhere else.

The class is paid off by passing to the double cover. Six servers carry the base
data repeated, a section exists upstairs in both classes, and the deck
transformation $\sigma(i) = i+3$ separates them. Full detail is in
\texttt{mobius.pdf}.

\section{The dictionary}

\begin{center}
\begin{tabular}{@{}>{\raggedright\arraybackslash}p{6.3cm}>{\raggedright\arraybackslash}p{7.5cm}@{}}
\toprule
\textbf{July} & \textbf{September} \\
\midrule
path directed container of $G$ & the ring of servers, one node per arc \\
edge function $f_e$ & the transition function on that overlap \\
the fold $F(p)$ & what \verb|Ask| accumulates in the backward maps \\
$F(p \oplus q) = F(q) \circ F(p)$ & the sign lives in the backward map \\
computable labelling & well-formed descent data \\
one class per vertex & a global section exists \\
Myhill--Nerode classes & the fibres of the minimal trivialising cover \\
class count per vertex & the index of that cover \\
memory needed at a vertex & the price of a centre \\
$F^n$ at refinement depth $n$ & \verb|Unroll| with fuel $n$ \\
\bottomrule
\end{tabular}
\end{center}

Two entries in that table deserve emphasis before the demonstration.

The fold law and the design constraint of the ring are the same equation. July
requires $F(p \oplus q) = F(q) \circ F(p)$ because that is what a directed
container's shift demands. September requires the sign to be applied on the way
back because descent data acts on sections and sections come back. Neither
document notices that the other has written down its own constraint.

And \verb|Unroll| is the refinement, one level at a time. The fibre returns
\verb|Exhausted| at every fuel, which says the initial algebra of the loop
composite is empty. The refinement's depth-$n$ table is the same object.

\section{The identification, demonstrated}

The refinement of \S\,``partition refinement'' was run on the path container of
the directed 3-cycle, with the edge functions being the M\"obius signs and then
the cylinder signs. The implementation groups paths by end vertex, then splits by
comparing continuations one level deeper, until a round splits nothing.

\begin{term}
mobius     signs=[1, 1, -1]  class count by step: [3, 6, 6, 6]
           stabilises at 6 classes;  reachable (vertex,value) pairs: 6
           [(0, -1), (0, 1), (1, -1), (1, 1), (2, -1), (2, 1)]
           quotient is a single cycle of length 6: [(0, -1), (1, -1), (2, -1), (0, 1), (1, 1), (2, 1)]

cylinder   signs=[1, 1, 1]  class count by step: [3, 3, 3, 3]
           stabilises at 3 classes;  reachable (vertex,value) pairs: 3
           [(0, 1), (1, 1), (2, 1)]
           quotient is a single cycle of length 3: [(0, 1), (1, 1), (2, 1)]
\end{term}

Read the M\"obius line against the program. The refinement splits once, from
three classes to six, and the quotient is a single 6-cycle. That is the ring
\texttt{./mobius-up.sh -\/-cover 2} brings up on ports 8811 to 8816. The values
attached to the six classes are $+1,+1,+1,-1,-1,-1$, which is the section
\texttt{./mobius cover 2} builds from node $0$ and installs one value per node.
The cylinder stabilises at three classes and its quotient is the base.

\begin{sectionbox}
\textbf{The identification.} The Myhill--Nerode quotient of the path container
is the cover that the cohomology class determines, and the classes are the
section. The minimal machine for a labelling has one state per reachable pair
$\langle \mathrm{end}(p), \ell(p) \rangle$, and its transition is
$(v,x) \mapsto (v+1,\, g_v \cdot x)$. That is the definition of the cover
associated to $\alpha : \pi_1 \to \mathbb{Z}/2$, and the transition is the
deck action.
\end{sectionbox}

The count is not a coincidence of $\mathbb{Z}/2$. Replacing the coefficients by
$\mathbb{Z}/m$ with one twist on the last edge gives a holonomy of order $m$,
and the class count follows it exactly.

\begin{term}
=== the minimal cover, for a 3-cycle with holonomy of order m ===
  Z/1: holonomy order  1   classes   3   = 3 x 1
  Z/2: holonomy order  2   classes   6   = 3 x 2
  Z/3: holonomy order  3   classes   9   = 3 x 3
  Z/4: holonomy order  4   classes  12   = 3 x 4
  Z/6: holonomy order  6   classes  18   = 3 x 6
\end{term}

Stated carefully, the class count at a vertex is the size of the orbit of the
seed value under the holonomy group. That equals the index of the smallest cover
on which this section becomes deck-invariant. When the seed is a fixed point the
orbit is a single element, the quotient is the base, and the labelling glued
already.

\section{Three diagnoses, three outcomes}

The July document ends its refinement section with a three-way diagnosis. The
program insists on distinguishing exactly three cases, and the README says why:
bad transition data and a non-zero class are different failures, and a
demonstration unable to tell them apart is worthless. The lists agree.

\begin{center}
\begin{tabular}{@{}>{\raggedright\arraybackslash}p{5.6cm}>{\raggedright\arraybackslash}p{4.1cm}>{\raggedright\arraybackslash}p{3.9cm}@{}}
\toprule
\textbf{July diagnosis} & \textbf{the federation} & \textbf{what the program says} \\
\midrule
one class per vertex, the future depends on the vertex alone
  & glues, a global section exists
  & cylinder, 2 sections \\
classes in bijection with the values, the labelling is computable
  & computable, and still no section
  & mobius band, 0 sections \\
two classes sharing a value, a collision
  & not descent data at all
  & the descent data is malformed \\
\bottomrule
\end{tabular}
\end{center}

The middle row is the one that sharpens both documents.

\begin{alarmbox}
\textbf{The M\"obius federation is computable.} Its labelling is a fold of edge
functions, which is exactly what \verb|Ask| accumulates. There is no collision
witness and nothing is uncomputable. What fails is only that the labelling is a
function of the \emph{path} and not of the \emph{endpoint}.

Descent asks for a labelling that factors through \verb|end|. The comonad
delivers one that factors through the path. The cohomology class measures the
gap, and by the July sentence quoted above, a non-zero class \emph{is} hidden
state. Here it is one bit of it.
\end{alarmbox}

The third row is a conjecture and is not demonstrated. Malformed transition data
should destroy the fold equation and so produce a collision witness, but the
argument is delicate and I have not run it. On the directed 3-cycle each node
applies its own sign when crossing outwards, so a cw-only path container still
has a well-defined edge function per edge and there is no collision. The
disagreement between neighbours is only visible once both directions are
positions, and the doubled graph needs its own refinement run. That is the first
thing to check.

\section{What the price of a centre actually is}

\texttt{sheaf-federation.pdf} says the centre is available at the price of
infinitely many copies of the federation. The identification shows that this is
wrong as stated, and shows what the right statement is.

The Myhill--Nerode quotient is the \emph{minimal} cover on which the section
becomes deck-invariant. Its index is the order of the holonomy. So with a finite
coefficient group and honest descent data, the price of a centre is always
finite, and it is exactly the order of the holonomy. For $\mathbb{Z}/2$ it is
two copies and six nodes, which the program pays in full.

The universal cover is a different object and a much larger bill. Consider the
figure-eight graph, one vertex and two independent loops, with coefficients in
$\mathbb{Z}/2$.

\begin{term}
=== figure-8: two independent loops, Z/2 coefficients ===
  minimal cover: 2 classes
  universal cover of the figure-8 is the infinite 4-regular tree
\end{term}

Two classes. The fundamental group is free of rank two and the universal cover
is an infinite tree, so passing to it would buy infinitely many copies in order
to kill a class that two copies already kill.

\begin{sectionbox}
\textbf{The correction.} Passing to the universal cover is sufficient and is
almost never necessary. What the class actually costs is the index of the
smallest cover that trivialises it, and the refinement computes that index
rather than bounding it. For finite coefficients the cost is finite. The phrase
``infinitely many copies'' describes one way of paying and not the price.
\end{sectionbox}

\section{Where ``almost computable'' actually lives}

That correction has a consequence that surprised me, and it relocates the
Sturmian material rather than applying it.

If the coefficient group is finite and the labelling is a fold, the holonomy has
finite order, the cover is finite, the refinement stabilises, and the labelling
is computable. Always. So no federation with finite local state and well-formed
descent data is ever almost computable. The whole Sturmian regime is unreachable
from inside descent theory.

There are exactly two ways out, and they are not variations of each other.

\textbf{Infinite coefficients.} Take $\mathbb{Z}$ rather than $\mathbb{Z}/m$.
The value set is infinite, so infinitely many classes appear at refinement depth
zero and the count never grows gradually. This is worse than Sturmian rather
than an instance of it, and the grading has little to say about it.

\textbf{A labelling that is not a fold.} Keep the value set finite and give up
the fold equation. Now the classes can be unbounded while the values are not,
and the splitting rate becomes the interesting invariant. This is the July
experiment. On the one-loop graph the refinement collapses to counting distinct
windows, and the Fibonacci word attains the minimum.

\begin{term}
=== a labelling that is NOT a fold: Fibonacci word on the one-loop graph ===
  c( 0) =   2   (n+2 = 2)
  c( 1) =   3   (n+2 = 3)
  c( 2) =   4   (n+2 = 4)
  c( 3) =   5   (n+2 = 5)
  c( 5) =   7   (n+2 = 7)
  c( 9) =  11   (n+2 = 11)
  c(13) =  15   (n+2 = 15)
  c(17) =  19   (n+2 = 19)
  c(21) =  23   (n+2 = 23)
  c(25) =  27   (n+2 = 27)
\end{term}

One new class per level, exactly, for as far as the word was generated.

\begin{sectionbox}
\textbf{The almost computable federations are the ones that are not federations.}
A labelling with a finite value set and an unbounded refinement is not generated
by any assignment of functions to edges. There is no cocycle, so there is no
class, and sheaf cohomology has nothing to say about it at all. The refinement
still grades it, and Morse--Hedlund still says that once it is infinite it grows
at least linearly.

So the July grading is not a refinement of $H^1$ on the region where $H^1$ is
defined. It is defined on a strictly larger region, and it agrees with $H^1$
where they overlap. That is a better result than the one I expected to find.
\end{sectionbox}

\section{What is solid and what is modelling}

\begin{alarmbox}
\textbf{Demonstrated in this note.} The refinement on the M\"obius ring gives six
classes, the quotient is the 6-cycle, and the classes carry the section
$+,+,+,-,-,-$. The class count is three times the order of the holonomy, for
$m = 1,2,3,4,6$. The figure-eight with $\mathbb{Z}/2$ coefficients has a
two-element minimal cover. The Fibonacci word gives $c(n) = n+2$ to depth
twenty-five.

\textbf{Solid, and neither mine nor in question.} Theorem~\ref{thm:july} and the
coKleisli collapse are July's. Myhill--Nerode, Morse--Hedlund and the
classification of covers by subgroups of $\pi_1$ are standard.

\textbf{Conjectured.} That malformed descent data is exactly the collision case.
That the identification of the Myhill--Nerode quotient with the minimal
trivialising cover holds for non-abelian coefficient groups, where $H^1$ is a
pointed set rather than a group and the orbit argument needs redoing.

\textbf{Open, and the reason to write this down.} Degree two. The triple overlap
of the three-arc cover is empty, so the cocycle condition is vacuous, and
nothing in the July document is about coherence either. Both constructions stop
at the same place for the same reason. A fourth arc, or a two-dimensional base,
is where either of them would first have something new to say.
\end{alarmbox}

\section{What to do next}

Three things, smallest first.

Run the refinement on the doubled graph, where both directions are positions,
and check whether malformed descent data produces a collision witness. That
settles the third row of the table in \S5 and it is an afternoon.

Replace $\mathbb{Z}/2$ by a non-abelian group. The smallest useful choice is
$S_3 \cong \mathrm{GL}_2(\mathbb{F}_2)$, where the classes are the conjugacy
classes and there are three of them rather than two. The prediction is that the
refinement returns three distinct minimal covers, of index one, two and three.
That is falsifiable, and it closes three limitations the federation README
currently admits to: the coefficients would stop being abelian, the two
directions round the ring would stop agreeing, and the forward map would acquire
real work.

Then degree two, which is the only one that is research rather than
implementation.

\end{document}
